\cotefile{15}
\scanpage{68}
$C$ catégorie de Deuring. On a alors la gr-catégorie $\underline{\mathrm{Equ}}(C)$, qui opère sur $C$, donc

\[
\boxed{\mathrm{Brd}(C)=\pi_0(\underline{\mathrm{Equ}})\quad
\deleted{\text{classes d'isom de }C}\ \text{opère sur Cl. d'isom des }C}
\]

Ceci dit, le quotient est l'ens des classes d'isomorphie des $A$-Algèbres $\unsure{o|\sigma}$ apparentées à $C$ (donc, dans le cas des \unsure{$K$|$\mathrm R$} ordres maximaux\ldots, \inserted{le quotient est formé, \unsure{$\exists$}} des classes de conjugaison de tous ordres maximaux, tandis que $\pi_0(\underline{\mathrm{Equ}})$ correspond \deleted{aux} \inserted{est le groupe des} classes d'isomorphie de $o$-bimodules $=$ \unsure{bi-$o$-réseaux} de $o$ mod aux \unsure{princ.|puiss.|provenant} de \deleted{\unsure{$K$|Pic}} $A^*$\unsure{\}|?}, \deleted{\unsure{Noter|Notons}} et $\mathrm{Iso}(C)$ est l'ens des classes d'isomorphie de $o$-modules : \unsure{droites|droits} ``inv.'' \deleted{de} ou encore des $o$-réseaux \deleted{\unsure{x}} à droite]. Dans le cas \inserted{des} \unsure{$A$ centraux|centraux|$o$ centraux} \deleted{\illegible{2}} \inserted{sur $A$} [on dit $C$ centrée sur $A$) on a de plus un foncteur \uline{\unsure{pleint} fidèle}

\[
\begin{aligned}
\underline{\underline{\mathrm{Pic}}}(A)&\longrightarrow\underline{\mathrm{Equ}}(C)\\
L&\longmapsto(X\longmapsto L\otimes_A X)
\end{aligned}
\]

d'où une suite exacte

\[
0\unsure{\mathrel{-}|\longrightarrow}\mathrm{Pic}(A)\longrightarrow
\deleted{\unsure{\mathrm{Br}}}\mathrm{Brdt}(C)\longrightarrow
\frac{\mathrm{Brdt}(C)}{\mathrm{Pic}(A)}\longrightarrow 0
\]

Dans le cas de Brauer (et ordres maximaux, $o$)

\[
\frac{\mathrm{Brd}\,C}{\mathrm{Pic}(A)}\simeq
\frac{\mathrm{Div}(C)}{\mathrm{Div}(A)}=
\underset{\substack{p\in\mathrm{Max}\,A\\p\ \text{ramifié en }C}}{\coprod}
\mathbb{Z}/\underset{\text{ramification en }p.}{e_p}\mathbb Z
\]
% archive sheet number: 67

\scanpage{69}
\mspage{- 30 -}

D'autre part, on a, en tenant compte de (26)

\[
\begin{aligned}
\frac{d}{dt}(tU_{kj}(\handwritten{x}_0+t(\handwritten{x}-\handwritten{x}_0);\handwritten{v}(t)))&=U_{kj}(\handwritten{x}_0+t(\handwritten{x}-\handwritten{x}_0);\handwritten{v}(t))+\\
&\quad+t\handwritten{\sum_i}(x_i-x_{i0})\frac{\handwritten{\partial}U_{kj}}{\handwritten{\partial} x_i}(\handwritten{x}_0+t(\handwritten{x}-\handwritten{x}_0);\handwritten{v}(t))\\
&\quad+t\handwritten{\sum_{i,h}}(x_i-x_{i0})\frac{\handwritten{\partial} U_{kj}}{\handwritten{\partial} y_h}(\handwritten{x}_0+t(\handwritten{x}-\handwritten{x}_0);\handwritten{v}(t))U_{ih}(\handwritten{x}_0+t(\handwritten{x}-\handwritten{x}_0);\handwritten{v}(t))
\end{aligned}
\]

d'où, en faisant intervenir (24), on voit que les fonctions

\[
w_{kj}(t)=\frac{\handwritten{\partial} v_j}{\handwritten{\partial} x_k}-tU_{kj}(\handwritten{x}_0+t(\handwritten{x}-\handwritten{x}_0);\handwritten{v}(t))\qquad
(1\leq j\leq q,\ 1\leq k\leq p)
\]

vérifient le système linéaire et homogène des $pq$ équations différentielles

\[
\tag*{(27)}\frac{dw_{kj}}{dt}=\handwritten{\sum_{i,h}}(x_i-x_{i0})
\frac{\handwritten{\partial} U_{ij}}{\handwritten{\partial} y_h}w_{kh}
\]

Or $w_{kj}(0)=0$, car pour $t=0$, $\handwritten{v}(0;\handwritten{x}_0,\handwritten{x},\handwritten{y}_0)=\handwritten{y}_0$, donc $\frac{\handwritten{\partial} v_j}{\handwritten{\partial} x_k}(0)=0$ ; comme $0$ est la seule solution de (27) qui s'annule pour $t=0$, on a identiquement $w_{kj}(t)=0$ pour $|t|<2$, et en particulier pour $t=1$ ; mais si on pose $\handwritten{u}(\handwritten{x},\handwritten{x}_0,\handwritten{y}_0)=\handwritten{v}(1;\handwritten{x}_0,\handwritten{x},\handwritten{y}_0)$, cela signifie que l'on a

\[
\frac{\handwritten{\partial}u_j}{\handwritten{\partial}x_i}=U_{ij}(\handwritten{x};\handwritten{u}(\handwritten{x}))\qquad
(1\leq i\leq p,\ 1\leq j\leq q)
\]

pour $\handwritten{x}\handwritten{\in}U$. De ces relations on déduit que si $\handwritten{z}(t)=\handwritten{u}(\handwritten{x}_0+t(\handwritten{x}-\handwritten{x}_0))$, on a

\[
\handwritten{z}'(t)=\underline{U}(\handwritten{x}_0+t(\handwritten{x}-\handwritten{x}_0);\handwritten{z}(t)).(\handwritten{x}-\handwritten{x}_0)
\]

pour $|t|<2$, et $\handwritten{z}(1)=\handwritten{v}(1)$ ; donc $\handwritten{z}(t)=\handwritten{v}(t)$\unsure{.|,} pour $|t|\leq 1$, et en particulier $\handwritten{u}(\handwritten{x}_0)=\handwritten{z}(0)=\handwritten{v}(0)=\handwritten{y}_0$ ; ce qui achève la démonstration.

On pourrait évidemment ici passer de la forme non homogène à la forme homogène du th. de Frobenius, avec l'algèbre extérieure, parler de la primitive (locale) d'une forme extérieure holomorphe, etc..; est-ce le lieu, ou,vaut-il mieux attendre les variétés analytiques pour le faire ?
% archive sheet number: 68

\scanpage{70}
Supposons $A$ \deleted{local} Dedekind \deleted{\unsure{$\ne$}} local \inserted{(cas Brauer)} de sorte que

\[
\begin{aligned}
\mathrm{Is}(C)&=1,\qquad \mathrm{Brd}(C)=\mathbb{Z}/\underset{\text{ramification}}{e}\mathbb{Z}\\
\mathrm{Pic}(A)&=0
\end{aligned}
\]

et on trouve

\[
\mathrm{Aut\,ext}(o)=\mathrm{Brd}\, C=\mathbb{Z}/e\mathbb{Z}
\]

Donc il y a obstruction locale à ce qu'un autom. de $o$ soit \unsure{néc.|nec.} intérieur, il y a $e$ classes d'autos de $o$, \ldots

Supposons \inserted{dans cas global} qu'un autom de $o$ soit localement intérieur (pour la top. \unsure{fppqc}, par exemple) \inserted{i.e. son image dans $\mathrm{Ram}(C\unsure{/\mathfrak p|/A})$ nulle}. \deleted{Alors les} \inserted{Il définit} donc une section de $\underline{o}^*/\mathrm{Gm}_A$, et \deleted{\unsure{elle est}} \inserted{l'obstruction} pour le relever est un élément de $\mathrm{Pic}(A)$, nul dans \deleted{\unsure{$\mathrm{Pic}(o)$}} $\mathrm{H}^1(A,o^*)$\ldots\ donnant un élément trivial de $\mathrm{Is}(C)$\ldots\ \deleted{\unsure{Donc}} ce qui se voyait aussi par les suites exactes précédentes

\begin{center}\rule{0.22\linewidth}{0.3pt}\end{center}

Cas des corps gauches de Deuring splittés en $p$ et $\infty$ : la suite exacte devient

\diagram{\begin{tikzcd}[ampersand replacement=\&,column sep=small]
0 \arrow[r,dash] \& \mathbb Z/2\mathbb Z\arrow[r] \& o^*\arrow[r] \& \mathrm{Aut}(o)\arrow[r] \& \mathbb Z/2\mathbb Z\arrow[r] \& \pi_0(C,x)\arrow[r] \& \pi_0(O,o)\arrow[r,dash] \&1
\end{tikzcd}}

\margin{engendré par $-1$}

\margin{$\Vert$ $\mathbb Z/2\mathbb Z$ sauf si $j=0$ ou $j=12^3$ (où on trouve $\mathbb Z/4\mathbb Z$ resp $\mathbb Z/6\mathbb Z$ si $p\ne 2,3$) \deleted{\unsure{et}}}

% slanted
\margin{on peut dire pour quels $p$ ces $j$ sont d'inv. de Hasse nul\ldots}

\margin{engendré par Frobenius}

\margin{épim ssi la courbe elliptique $X$ d'anneau $o$ se descend à $\mathbb F_p$, i.e. \unsure{ssi} $j\in\mathbb F_p$}

\margin{classes d'iso. des classes de courbes elliptiques de Hasse nul l'opération de ``Frobenius'' dessus est $E\longmapsto E^{(p)}$}

\margin{classes de conjugaison possibles d'ordres maximaux}

Question : $o^*$ est-il tjrs égal
% archive sheet number: 69

\scanpage{71}
\mspage{- 25 -}

\[
\left.\begin{aligned}
H_i(0;\mu_1,\ldots,\mu_q)&=0\\
mH_i(\unsure{n|m};\mu_1,\ldots,\mu_q)&=\handwritten{\sum'_{n,(r_j),(\lambda_\ell)}}G_i(n;r_1,\ldots,r_p;\handwritten{\lambda}_1,\ldots,\handwritten{\lambda}_q)\mathbin{.}\\
&\quad\mathbin{.}\handwritten{\prod_{j=1}^{p}\prod_{k=1}^{r_j}}H_i(s_{kj};\handwritten{\nu}_{kj1},\ldots,\handwritten{\nu}_{kjq})
\end{aligned}\right\}
\quad\text{pour }m>0.\tag*{(17)}
\]

la sommation étant étendue aux indices satisfaisant aux mêmes conditions (16) que dans (15). La comparaison de (15) et (17) montre par récurrence sur le degré total que l'on a

\[
\left|h_i(m;\mu_1,\ldots,\mu_{\unsure{q|c}})\right|\leq H_i(m;\mu_1,\ldots,\mu_q)
\]

pour tout système d'indices ; en d'autres termes, les séries $(H_i(m;\mu_1,\ldots,\mu_q)t^mw_1^{\mu_1}\ldots w_q^{\mu_q})$ sont des \uline{majorantes} pour les séries $(h_i(m;\mu_1,\ldots,\mu_q)x^mz_1^{\mu_1}\ldots z_q^{\mu_q})$ respectivement $(1\leq i\leq p)$. Tout revient encore à prouver qu'on peut choisir les $G_i(n;r_1,\ldots,r_p;\handwritten{\lambda}_1,\ldots,\handwritten{\lambda}_q)$ de sorte que les séries majorantes $(H_i(m;\mu_1,\ldots,\mu_q)t^mw_1^{\mu_1}\ldots w_q^{\mu_q})$ soient sommables dans un voisinage de $0$.

Or, la série entière dont la somme est fonction des $p+q+1$ variables réelles $t$, $v_i$ $(1\leq i\leq p)$ et $w_{\handwritten{\ell}}$ $(1\leq\handwritten{\ell}\leq q)$

\[
\frac{M}{(1-\alpha t)(1-\alpha\handwritten{\sum_{i=1}^{p}}v_i)(1-\alpha\handwritten{\sum_{\ell=1}^{q}}w_{\handwritten{\ell}})}
\]

est une \deleted{am}\handwritten{ma}jorante de chacune des séries (14) dès que $M$ est assez grand et $\alpha$ assez petit. Tout revient à prouver qu'il existe $p$ séries entières sa ns terme indépendant de $t$, qui satisfont aux équations

\[
\frac{dv}{dt}\,k=\frac{M}{(1-\alpha t)(1-\alpha\handwritten{\sum_{i=1}^{p}}v_i)(1-\alpha\handwritten{\sum_{\ell=1}^{q}}w_{\handwritten{\ell}})}
\qquad(1\leq k\leq p)\tag*{(18)}
\]

On déduit aussitôt de (18) que si ce système admet une solution, tous les $v_k$ sont égaux à $v_1=v$ satisfaisant à l'équation différentielle
% archive sheet number: 70

\scanpage{72}
Soit \deleted{\unsure{$(C_x)$} \illegible{1}} $O$ la catégorie des $A$-Algèbres apparentées \deleted{\unsure{autrement}} à une $A$-Algèbre déterminée, \deleted{et pour tout \unsure{l'un deux} $o\in O$, \unsure{donc} supposons qu'il existe une fonction $S$ de $A$\textsuperscript{reg} \unsure{(cas}} Soit $K$ l'anneau total des fractions de $A$, \inserted{d'où $A\subset K$,} et supposons $K$ $1/2$ local.

Donc \deleted{\unsure{tt}} pour $o\in O$, $o'\in O'$, tt bimodule $L$/sur $o,o'$ loc. libre de rang 1 \inserted{de chaque côté} loc. sur $S=\mathrm{Spec}\,A$ est libre \inserted{séparément} sur $o_K$ \deleted{\unsure{$=$|$\simeq$}} $E$ et sur $o'_K$, d'où on conclut que
% slanted
\margin{$E$ à isom. près \unsure{ne dépend|on dispose|un défini} \unsure{pas de|par} $o$, seulement de $O$}
a) $o'_K\simeq E$ \deleted{\unsure{plongé dans} $E$} \inserted{donc $o'$ se plonge dans $E$}, et b) \inserted{\uuline{si tout autom de $E$ est intérieur}} tt $L$ comme dessus se plonge \deleted{\unsure{$B$.}} dans $E$ \inserted{avec la structure \unsure{d'g.|d.g.}} comme un $(o,o')$-réseau,

La sous $2$-catégorie \inserted{pleine \deleted{de} (Cat)} dont les objets sont \deleted{les \unsure{$A$} $O$, les} \inserted{les cat. de} Deuring de type $o$, \deleted{est} qui est équivalente : celle $\underline{\underline B}_1$ dont les objets sont les objets de $O$, $\underline{\mathrm{Hom}}(o,o')=$ \inserted{catégorie des} bimodules \unsure{sous|sur} $o\unsure{^{*}|^{}},o'$ qui sont loc \deleted{\unsure{sur $S$}} \inserted{libres de} rang 1 de chaque côté loc. sur $S$, et la loi de composit\unsure{ion|.} étant

\[
\begin{aligned}
&\deleted{\unsure{\mathrm{de}\ o}}\quad
\deleted{{}_oM_{o'},{}_{o''}\unsure{L}_{o'}\longrightarrow M\otimes_{o'}\unsure{L}}\\
&{}_oL_{o'},M_{o',o''}\longrightarrow L\otimes_{o'}M,
\end{aligned}
\]

$\underline{\underline B}_2$ est aussi équivalente au \uline{$2$-groupoïde de Brandt} dont les objets sont les sous-anneaux $o$ \deleted{apparentés} \inserted{de $E$} $\in O$, \deleted{\unsure{les}}

\[
\begin{aligned}
\underline{\mathrm{Hom}}(o,o')={}&\deleted{\text{catégorie des}}\inserted{\text{ensemble des}}\
 (o,o')\text{-réseaux }L\\
&\text{de }E\text{ tels que } \unsure{\text{ils sont}|\text{elles sont}}\\
&\text{des deux côtés loc. lib. de rg}\\
&1\text{ loc sur }S\\
\mathrm{Hom}(L,M)={}&\text{ens des }\lambda\in E^{\unsure{\leftrightarrow|*|\natural}}\text{ tels que }\lambda L=M.
\end{aligned}
\]

\margin{si $L,M$ sont des $(o,o')$ biréseaux ``spéciaux''}

\margin{centre de $E$}

On notera que les $2$-foncteurs $\underline{\underline B}_2\longrightarrow\underline{\underline B}_1\longrightarrow\underline{\underline B}\unsure{\overset{=}{\leftarrow}|=}\underline{\underline B}_{\unsure{o|0}}$ n'ont pas de quasi-inverses ``naturels'' !
% archive sheet number: 71

\scanpage{73}
\mspage{- 18 -}

\[
f_i(x_1,..,x_p\,;\,v_1(x_1,..,x_p),..,v_q(x_1,..,x_p))=0
\quad(1\leq i\leq q)\tag*{(2)}
\]

\uline{dans un voisinage de} $(a_1,..,a_p)$, \uline{on a} $u_i=v_i$ $(1\leq i\leq q)$ \uline{dans un voisinage de} $(a_1,..,a_p)$.

Par tra\handwritten{\unsure{-|{}}}nslation, on peut se ramener au cas où $a_j=0$ pour $1\leq j\leq p$ et $b_i=0$ pour $1\leq i\leq q$, . Posons pour $1\leq i\leq q$

\[
f_i(x_1,..,x_p;y_1,..,y_q)=f_{i0}+
\handwritten{\sum_{j=1}^{q}}f_{ij}y_j+
\handwritten{\sum_{(n_j)}}f_{\unsure{1|i};n_1..n_q}y_1^{n_1}\ldots y_q^{n_q}\tag*{(3)}
\]

les $f_{i0}$, $f_{ij}$ et $f_{i;n_1..n_q}$ étant des séries entières par rapport aux $x_k$ $(1\leq k\leq p)$ sommables dans un polycylindre ouvert $P$ de centre l'origine dans $K^p$, $f_{i0}$ n'ayant pas de terme constant ; en outre, $f_i$ est sommable dans $P\times Q$, où $Q$ est un polycylindre de centre l'origine dans $K^q$, et dans la seconde sommation des seconds membres de (3), on a toujours $\handwritten{\sum_{j=1}^{q}}n_j\geq 2$. Par hypothèse, la matrice $\underline F=(f_{ij})$ est holomorphe dans $P$ et inversible à l'origine donc elle est inversible dans un polycylindre ouvert contenu dans $P$, et son inverse $\underline G=(g_{ij})$ est holomorphe dans ce polycylindre ; nous supposerons donc chacune des $g_{ij}$ développées en série entière sommable dans un polycylindre ouvert $P_1\handwritten{\subset} P$ de centre l'origine ; si on pose alors $g_i=\handwritten{\sum_{j=1}^{q}}g_{ij}f_j$, on aura

\[
g_i(x_1,..,x_p;y_1,..,y_q)=-h_{i0}+y_i-
\handwritten{\sum_{(n_j)}}h_{i;n_1..n_q}y_1^{n_1}\ldots y_q^{n_q}\tag*{(4)}
\]

où $h_{i0}$ et les $h_{i;n_1..n_q}$ sont des séries entières par rapport aux $x_k$, sommables dans $P_1$, $h_{i0}$ n'ayant pas de terme constant, et la somme étant étendue aux $(n_j)$ tels que $\handwritten{\sum_{j=1}^{q}} n_j\geq 2$. La première partie du théorème revient à prouver qu'il existe $q$ séries entières sans terme constant $u_i$ $(1\leq i\leq q)$ par rapport aux $x_k$, sommables dans un voisinage de l'origine, et telles que l'on ait

\[
u_i=h_{i0}+\handwritten{\sum_{(n_j)}}h_{i;n_1..n_q}u_1^{n_1}\ldots u_q^{n_q}
\qquad\text{pour }\ 1\leq i\leq q\ .\quad .\tag*{(5)}
\]
% archive sheet number: 72

\scanpage{74}
Supposons maintenant la condition suivante satisfaite (en plus de : \uline{tout autom de} $E$ \uline{est intérieur}).

(C) Pour tout $o\subset E$, $o\in O$, la loi de composition de $\mathrm{Hom}_{\underline B_1}(o,o)=$ ensemble des \inserted{$(o,o)$-}biréseaux spéciaux \deleted{\unsure{sur $(o,o)$}} \deleted{est} $(L,M\longmapsto L.M)$ est commutative :

\[
LM=ML.
\]

% slanted
\margin{Il suffit (et il faut si $A$ est \unsure{Dedek.} de dim 1) que ce soit vrai sur les localisés \deleted{\unsure{en les}} idéaux maximaux (ou leurs complétés, si $A$ noeth.). Dans le cas local, \deleted{\unsure{il suffit}} c'est équiv. : les relations $LM=ML$ sont vraies pour $L,M$ des idéaux bilatères biprincipaux de $o$}

$C\in\mathrm{Ob}_0\,\underline{\underline B}_0$ une catégorie de Deuring, et $X\in\mathrm{ob}\,C$, de sorte que $o\unsure{\simeq|=}\mathrm{End}_C(X)\in O$ ; on peut \deleted{supposer} \inserted{choisir} $E\simeq o_K$ \inserted{(et ``identifier'')}. On va définir sur la gr-catégorie $\underline{\unsure{\mathrm{EquAut}|\mathrm{Equ}|\mathrm{Eqv}}}(C)=\underline{\mathrm{Hom}}_{\underline{\underline B}_0}(C,C)$ une \deleted{loi de} contrainte de commutativité, \deleted{\unsure{définit}} \inserted{en fait une \uline{cat. de Picard stricte}} à priori du choix de $X$. En effet, on a un gr-isomorphisme

\[
\underline{\mathrm{Hom}}_{\underline{\underline B}_0}(C,C)
\overset{\alpha_X}{\simeq}
\underline{\mathrm{Hom}}_{\underline{\underline B}_1}(C,C)
=\text{catégorie des }\deleted{\text{\unsure{bimodules} spéciaux sous }(o,o)}
\inserted{\text{$(o,o)$-réseaux}}\text{ spéciaux}
\]

cette dernière catégorie étant \inserted{cat. de Picard} strictmt comm. (en fait, les contraintes d'associativité \uline{et} de commutativité y sont des identités). On peut donc définir une contrainte de comm. stricte sur $\underline{\mathrm{Equ}}(C)$ par transport de structure. On peut l'expliciter ainsi : si $L,M$ sont des $(o,o)$-bimodules spéciaux, on définit

\[
c:L\otimes_o M\overset{\sim}{\longrightarrow}M\otimes_o L
\]

en \deleted{\unsure{défini}} utilisant des isom de $(o,o)$ bimodules

\[
\begin{aligned}L&\simeq L'\subset E\\M&\simeq M'\subset E\end{aligned}
\]

avec $L',M'$ des $(o,o)$-réseaux, et notant que

\[
\begin{aligned}
L\otimes M&\simeq L'M'.\\
&\phantom{\simeq L'}\shortparallel\\[-.5ex]
M\otimes L&\simeq M'L'
\end{aligned}
\]
% archive sheet number: 73

\scanpage{75}
\mspage{- 19 -}

On notera que tout système de $q$ séries entières $u_i$ sans terme constant $(1\leq i\leq q)$ sommables dans un voisinage de l'origine, est substituable da ns chacune des séries entières $\handwritten{\sum_{(n_j)}}h_{i;n_1..n_q}y_1^{n_1}..y_q^{n_q}$ (cor. du th.2) ; si le problème a une solution, les séries entières (par rapport aux $x_i$) des deux membres de (5) \unsure{sort|sont} donc sommables au voisinage de l'origine, et par suite doivent a\unsure{-|{}}voir leurs coefficients \uline{identiques}.

Supposons donc le problème résolu, et désignons par $c_i(n_1,..,{}_{q};r_1,..,r_p)$ le \unsure{cœfficient|coefficient} de $x_1^{r_1}\ldots x_p^{r_p}$ dans la série $h_{i;n_1..n_q}$, par $c_i(r_1,..,r_p)$ le coefficient du même terme dans $h_{i0}$ par $d_i(m_1,..,m_p)$ le coefficient de $x_1^{m_1}..x_p^{m_p}$ dans $u_i$. L'identification des séries des deux membres de (5) donne les conditions

\[
\left\{\begin{aligned}
d_i(0,0,\ldots,0)&=0\\
d_i(m_1,..,m_p)&=c_i(m_1,..,m_p)+\\
&\quad+\handwritten{\sum_{(n_j)(r_i)}}(c_i(n_1,..,n_q;r_1,..,r_p)
\handwritten{\prod_{j=1}^q\prod_{k=1}^{n_j}}d_j(t_{kj1},t_{kj2},\ldots,t_{kjp}))
\end{aligned}\right.\tag*{(6)}
\]

la sommation du second membre étant étendue à tous les systèmes d'entiers $\geq 0$, $n_j,r_i$ et $t_{kji}$ satisfaisant aux conditions $\handwritten{\sum_{i=1}^{p}}t_{kji}\geq 1$, $\handwritten{\sum_{j=1}^q}n_j\geq 2$ et

\[
m_i=r_i+\handwritten{\sum_{j=1}^{q}\sum_{k=1}^{n_j}}t_{kji}\qquad(1\leq i\leq p)\ .\tag*{(7)}
\]

On tire de ces conditions que

\[
\handwritten{\sum_{i=1}^{p}}m_i=\handwritten{\sum_{i=1}^{p}}r_i+
\handwritten{\sum_{j=1}^{q}\sum_{k=1}^{n_j}\sum_{i=1}^{p}}t_{kji}
\]

Si on tient compte de ce que $\handwritten{\sum_{i=1}^p}t_{kji}\geq 1$ et $\handwritten{\sum_{j=1}^{q}}n_j\geq 2$, on voit que pour tout système $(t_{kji})_{1\leq i\leq p}$ qui figure au second membre de (7), on a $\handwritten{\sum_{i=1}^p}t_{kji}<\handwritten{\sum_{i=1}^{p}}m_i$, en d'autres termes, les relations (6) déterminent les coefficients des termes de degré \uline{total} $s$ dans $u_i$ en fonction de ceux des termes de degré total $<s$ dans les $u_j$ $(1\leq j\leq q)$ ; par récurrence sur $s$, on voit donc que les équations (6) déterminent les coefficients
% archive sheet number: 74

\scanpage{76}
Il est immédiat que \deleted{cette} \deleted{\unsure{loi}} l'isom ne dépend pas du choix des \unsure{plongmt|plongemt} $L\hookrightarrow E$, $M\hookrightarrow E$.

\uline{Proposition} La contrainte de commutativité précédente ne dépend pas du choix de $X\in\mathrm{Ob}\,C$.

[\deleted{Ceci dit, en effet \unsure{un objet|un autre} $X'$ de} Il est clair qu'elle ne change pas si on remplace $X$ par un objet isomorphe [il faudrait une petite démonstration : dans une gr-catégorie $B$, les \deleted{foncteurs} \unsure{auto-intérieurs|automorphismes|autoéquivalences} \deleted{\unsure{opér|qui}} transforment une \deleted{\unsure{donnée}} \inserted{contrainte} de comm. (stricte ou non) en elle-m\unsure{ê|ême}\ldots ]]

\margin{$\pi_0$}

\deleted{On peut \unsure{d'ailleurs partir de}} Soit $X'$ un autre choix, défini par un $P\in\mathbb P(o/A)$. On peut supposer $P\subset E$, soit $o'=\mathrm{End}_{\unsure{o|{}}}(P)$, de sorte que $o'$ est maintenant plongé dans $E$ (c'est l'ens des \deleted{\unsure{Auto}} $\lambda\in E$ tels que $\lambda P\subset P$) Ainsi $P$ est un \deleted{birésea} $(o',o)$-réseau. On a un diagramme comm.

\diagram{\begin{tikzcd}[ampersand replacement=\&,column sep=small]
\& \underline{\mathrm{Equ}}(C,C)\arrow[dl,"\alpha_X"']\arrow[dr] \& \\
\substack{\text{Catégorie des}\\(o,o)\text{ réseaux}\\\text{spéciaux}}\arrow[rr,"{L\mapsto P\deleted{\unsure{\otimes}}LP^{-1}}"] \& \&
\substack{\text{Catégorie des}\\(o',o')\text{-réseaux}\\\text{spéciaux}}
\end{tikzcd}}

Il suffit donc de prouver que $L\longmapsto P.L$ \deleted{\unsure{transforme}} \inserted{est comp} \deleted{\unsure{les deux}} \inserted{avec les} contraintes de commutativité

\diagram{\begin{tikzcd}[ampersand replacement=\&,column sep=small]
LM\arrow[d,equal,"{c_X=1}"',"{\deleted{\unsure{1=c}}}"] \& \text{transformé par }P\text{ en} \&
P(LM)P^{-1}\arrow[d,equal,"{P(c_X)=1}"] \& {}=(PLP^{-1})(PMP^{-1})\arrow[d,equal,"{c_{X'}=1}"]\\
ML \& \& P(M.L)P^{-1} \& {}=(\deleted{\unsure{M}}PMP^{-1})(PLP^{-1})
\end{tikzcd}}

OK.
% archive sheet number: 75

\scanpage{77}
\mspage{- 17 -}

$|z_i|\leq r_i$ ; soit $M$ la borne supérieure des normes des termes de la série pour $|z_i|=r_i$ $(1\leq i\leq p)$ ; on a donc

\[
\left\|\handwritten{\unsure{c|\mathbb C|C}}_{n_1n_2..n_p}\right\|\leq M/r_1^{n_1}..r_p^{n_p}
\]

ce qui montre que la série entière $(M\,(\frac{x_1}{r_1})^{\handwritten{n}_1}\ldots(\frac{x_p}{r_p})^{\handwritten{n}_p})$ est une majorante de la série donnée ; on observera que la somme de cette majorante est $M/(1-\frac{x_1}{r_1})(1-\frac{x_2}{r_2})\ldots(1-\frac{x_p}{r_p})$ pour $|x_i|<r_i$ $(1\leq i\leq p)$. Une autre majorante est la série obtenue par substitution de $\handwritten{\sum_{i=1}^p}\frac{x_i}{r_i}$ à $u$ dans la série géométrique $(Mu^n)$ ; en effet, il est immédiat que le coefficient de $x_1^{n_1}..x_p^{n_p}$ dans cette série est le produit d'un entier $\geq1$ et de $M/r_1^{n_1}..r_p^{n_p}$ ; cette nouvelle majorante a pour somme $M/(1-\handwritten{\sum_{i=1}^p}\frac{x_i}{r_i})$ dans un voisinage assez petit de l'origine.

\uline{2}. \uline{Résolution des équations analytiques}.

THEOREME 1 .- \uline{Soient} $f_i(x_1,..,x_p\,;\,y_1,y_2,\ldots,y_p)$ $(1\leq i\leq q)$ $q$ \uline{fonctions holomorphes scalaires définies dans une partie ouverte} $A$ \uline{de} $K^{p+q}$. \uline{Soit} $(a_1,\ldots,a_p\,;\,b_1,..,b_q)$ \uline{un point de} $A$ \uline{tel que} $f_i(a_1,..a_p\,;\,b_1,..,b_q)=0$ \uline{pour} $1\leq i\leq q$, \uline{et que le déterminant fonctionnel} $\handwritten{\Delta}=\det(\handwritten{\partial} f_i/\handwritten{\partial} y_j)$ \uline{soit} $\handwritten{\ne}0$ \uline{en ce point}. \uline{Alors il existe un voisinage} $V$ \uline{de} $(a_1,..,a_p)$ \uline{dans} $K^p$ \uline{et un système de} $q$ \uline{fonctions scalaires} $u_i(x_1,..,x_p)$ $(1\leq i\leq q)$ \uline{holomorphes dans} $V$, \uline{satisfaisant aux conditions} $u_i(a_1,..,a_p)=b_i$ $(1\leq i\leq q)$ \uline{et}

\[
f_i(x_1,\ldots,x_p\,;\,u_1(x_1,..,x_p),..,u_q(x_1,..,x_p))=0\qquad(1\leq i\leq q)\tag*{(1)}
\]

\uline{en tout point de} $V$. \uline{En outre}, \uline{si} $v_i$ $(1\leq i\leq q)$ \uline{est un système de} $q$ \uline{fonctions continues au point} $(a_1,\ldots,a_p)$, \uline{définies dans un voisinage de ce point}, \uline{et satisfaisant aux conditions} $v_i(a_1,..,a_p)=b_i$ $(1\leq i\leq q)$ \uline{et}

% handwritten notes in the opposite direction to the typescript
\handwritten{$\deleted{\unsure{\mathbb Z_2|2}}\ \mathbb Z_p($}

\handwritten{$Q(i)^{\deleted{\unsure{*}}}\qquad T$}

\handwritten{$\unsure{B^o_C|B^0_C|B^o_c|\mathcal{B}_c}(\mathbb Z)$}

\handwritten{$-1$ \inserted{pas un} carré mod $p$\qquad $p\equiv-1\qquad(4)$}

\handwritten{$p\qquad\left[\genfrac{}{}{0pt}{}{-1}{p}\right]=-1\qquad p$}

\handwritten{\unsure{TORS|TORs}(}
% archive sheet number: 76

\scanpage{78}
\uline{Conclusion} Les catégories $\underline{\mathrm{Equ}}(C)$, $C\in\mathrm{ob}_0\underline{\underline B}_0$, sont \deleted{des gr-catégories munies de str} \inserted{munies canoniquement d'une structure de} catégorie de Picard \uline{stricte\deleted{s}} \deleted{\uline{\unsure{commutatives}}}. \deleted{Ils sont} Ces cat. de Picard sont \uline{canoniquement} équivalentes entre elles, et peuvent donc s'identifier à une m\textsuperscript{e} catégorie de Picard $B$ (catégorie de Picard de \uline{Brandt}). Donc $\underline B$ opère sur chaque $C$, et pour $C^{\unsure{*|{}}},C'\in\mathrm{ob}_0\underline{\underline B}_0$, la catégorie \deleted{Hom} $\underline{\mathrm{Equ}}(C,C')$ est canoniquement munie d'une structure de tors-catégorie sous $\underline B$. De façon précise,

\[
\begin{aligned}
C'&\longmapsto\underline{\mathrm{Equ}}(C,C')\\
\underline{\underline B}_0&\longrightarrow \text{2-catégorie des tors-catégories}\\
&\hspace{5em}\text{sous }\underline B
\end{aligned}
\]

est une \unsure{équivalence|équivalente} de $2$-catégories.

Si on a un chgt de base

\diagram{\begin{tikzcd}[ampersand replacement=\&,row sep=small,column sep=small]
A\arrow[r]\arrow[d,hook,dash] \& A'\arrow[d,hook,dash]\\
K\arrow[r,hook,dash] \& K'
\end{tikzcd}}

, et si les $C_{A'}$ satisfont les m\textsuperscript{es} conditions que les $C$, alors on a un hom de \deleted{gr} catégories de Picard

\[
\underline B\longrightarrow B'
\]

et l'hom \deleted{\unsure{de} \illegible{2}}

\[
C\longrightarrow C_{A'}
\]

est compatible avec les opérations de $B,B'$\ldots
% archive sheet number: 77

\scanpage{79}
\mspage{- 16 -}

ce qui prouve que la série considérée est absolument convergente pour $\|\handwritten{u}\|<1/5a$. Le fait que $(\handwritten{h}(\handwritten{u}))^2=\handwritten{e}+\handwritten{u}$ résulte de ce que la formule $(1-\handwritten{\sum_{n=1}^{\infty}}(\genfrac{}{}{0pt}{}{2n}{n})x^n)^2=1-4x$ est vraie lorsque $E=K=\handwritten{\mathbb R}$ (\uline{Fonct.var.réelle}, chap.III, \S\,2) (cf.n\textsuperscript{o}2, remarque suivant la prop.2).

Reste à montrer que la fonction $\handwritten{g}$ est la seule fonction continue et égale à $\handwritten{e}$ à l'origine, dont le carré soit $\handwritten{f}$. Or, la relation $\handwritten{g}^2-\handwritten{g}_1^2=0$ s'écrit $(\handwritten{g}-\handwritten{g}_1)(\handwritten{g}+\handwritten{g}_1)=0$ puisque $E$ est commutative ; mais comme $\handwritten{g}(0)+\handwritten{g}_1(0)=2\handwritten{e}$, $\handwritten{g}+\handwritten{g}_1$ est inversible dans un voisinage de l'origine, donc on a $\handwritten{g}-\handwritten{g}_1=0$ dans ce voisinage, ce qui achève la démonstration.

\S\,2. La méthode des majorantes.

\uline{1}. \uline{Majorantes d'une série entière}.

DEFINITION 1.- \uline{Etant donnée une série entière} $(\handwritten{\unsure{c|C}}_{n_1n_2..n_p}z_1^{n_1}..z_p^{n_p})$ \uline{à coefficients dans} $E$, \uline{par rapport aux variables} $z_i\in K$ $(1\leq i\leq p)$, \uline{on dit qu'une série entière} $(C_{n_1n_2..n_p}x_1^{n_1}..x_p^{n_p})$ \uline{à coefficients positi\unsure{fs}} \uline{par rapport aux variables réelles} $x_i$ $(1\leq i\leq p)$, \uline{est une majorante de l\unsure{a} série entière donnée,} \uline{si pour tout système d'exposants} $(n_1,n_2,\ldots,n_p)$, \uline{on a} $\|\handwritten{c}_{n_1n_2..n_p}\|\leq C_{n_1n_2..n_p}$.

Il est clair que la somme (resp. le produit) de majorantes de deux séries entières est une majorante de la somme (resp. du produit) de ces deux séries.

Si une majorante d'une série entière est sommable pour $x_1=r_1>0,\ldots,x_p=r_p>0$, il est clair que la série entière donnée est absolument somma\unsure{-|{}}ble dans le polycylindre fermé défini par $|z_i|\leq r_i$ $(1\leq i\leq p)$. Inversement, supposons qu'une série entière $(\handwritten{c}_{n_1..n_p}z_1^{n_1}\ldots z_p^{n_p})$ soit absolument sommable dans le polycylindre fermé
% archive sheet number: 78

\scanpage{84}
\margin{$\unsure{\epsilon|E|c}$}

$\mathcal M$ catégorie motivique \inserted{$1/2$ simple}, de gerbe $\mathscr G$. Supposons que le lien $\mathcal L$ de $\mathscr G$ soit \uline{commutatif}, donc correspond à un groupe algébrique \inserted{comm de type multiplicatif} $G$ sur $\mathbb Q$. Ainsi, $G^0$ est un tore. Soit $E$ un motif ``fidèle''. \deleted{et soit $F$ un foncteur} \inserted{Alors $G$ opère} sur $E$, considéré comme \deleted{objet} \inserted{section} de la catégorie fibrée $\mathrm{HOM}(\mathscr G,\underline{\underline{\unsure{\mathrm{Proj}|\mathrm{Rep}|\mathrm{Prj}}}})$, et en particulier $G(\mathbb Q)$ opère sur $E$ donc $G(\mathbb Q)\overset{\varphi}{\hookrightarrow}\mathrm{End}_{\mathcal M}(E)$. D'ailleurs,

\deleted{les \unsure{élémts} de $G(\mathbb Q)$}

\deleted{Soit $U=\mathrm{Ker}(G^0\longrightarrow G_m)$, alors $U$ opère}

\deleted{\unsure{La}} Choisissons une forme de polarisation sur $E$, d'où une structure d'algèbre de Riemann sur $A=\mathrm{End}_{\mathcal M}(E)$, et on trouve que pour $u\in G(\mathbb Q)$, $\varphi(u)\varphi(u)^{\deleted{\unsure{*}}'}=\varepsilon(u)1\deleted{\unsure{*|E}}$.

[En particulier, si $u\in\inserted{U=}\mathrm{Ker}\,\varepsilon$, on aura $\varphi(u)\varphi(u)'=\mathrm{id}$.]. Donc l'algèbre engendrée par les $\varphi(u)$ dans $A$ est auto-adjointe, et $U(\mathbb Q)$ apparaît comme un ss-groupe du ``groupe unitaire'' de $A$, formé des $u$ tels que $uu'=1$. \deleted{Plus précisément} Or ce \inserted{dernier} groupe est \inserted{(formé sur $\unsure{\mathrm{Vect}.|\mathrm{Vert}.|\mathrm{vect}}/\mathbb Q$) et} un \uline{tore} \uline{compact}. On trouve que $U^0$ est un \uline{tore compact}.
% archive sheet number: 83

\scanpage{85}
où $p_{X'}$ désigne le morphisme structural $X'\quad X$. D'autre part, pour donner un sens à 4.1. (ii) et justifier \deleted{\unsure{sa|la}}la démonstration indiquée de 4.1. , il y a lieu d'intorduire un univers \quad tel que \qquad \unsure{,|;} et considérer les faisceaux intervenant dans les formules \deleted{\unsure{5}}4.1. (i\unsure{i)} comme des faisceaux à valeurs dans \quad -(Ens).

\uline{Proposition} 4L1. (i) Le foncteur $u^x$ est pleinement fidèle, donc pour tout faisceau $G$ sur $X_{\text{ét}}$, l'homomorphisme canonique $G\qquad u_x u^x(G)$ est un isomorphisme.

(ii) Soient $G$ (resp. $F$) un faisceau abélien sur $X_{\text{ét}}$ (resp. $(\mathrm{Sch})_{/X}$). Alors les homomorphismes canoniques (définis par exemple comme edge-homomorphismes de la suite \unsure{spectrahe|spectrale} de Leray V 5.5.1.) suivants sont\deleted{\unsure{s|x}} des isomorphismes:

\[
\begin{aligned}
H^{\deleted{\unsure{x|{}}}*}(X_{\text{ét}},F\,u)&\simeq H^{\deleted{x}*}((\mathrm{Sch})_{/X},F)\\
H^*(X_{\text{ét}},G)&\simeq H^{\deleted{x}*}((\mathrm{Sch})_{/X},u^{\deleted{x}*}(G))\ .
\end{aligned}
\]

En effet, le foncteur $u$ est un comorphisme (VI 7.3.6.) et on peut alors appliquer VI 7.3.7. et VI 7.3.9.
% archive sheet number: 84

\scanpage{86}
On peut préciser, en introduisant l'Algèbre engendrée par \deleted{$A$} $G(\mathbb Q)$ dans $A$, soit $B$.

On sait que $B=\prod L_i$, chaque $L_i$ ext. quadratique d'un corps de nbres $K_i$ tot. réel. ($K_i$ déterminé uniqmt \deleted{par la condition} \inserted{dans $L_i$} comme l'ens des élémts tot. réels de $L_i$), l'involution est donnée sur $L_i$ par la conjugaison p.r. à $K_i$ (i.e. la conjugaison complexe ordinaire, si $L_i\subset\mathbb C$), la trace est proportionnelle, sur chaque $L_i$, à la trace ordinaire.

Ainsi $B^*$ n'est autre que $H(\mathbb Q)$, où

\[
H=\prod_i\left(\prod_{L_i/\mathbb Q}\mathrm{Gm}_{L_i}\right)
\]

et le groupe unitaire $B^{*u}$ n'est autre que le \deleted{groupe} des pts rationnels de

\[
H^u=\prod_i\left(\prod_{K_i/\mathbb Q}\Gamma_i\right)
\]

où $\Gamma_i$ est la forme tordue de $\mathrm{Gm}$ sur $K_i$ définie par l'extension quadratique $L_i$ de $K_i$.

\deleted{[} \deleted{Ainsi} \unsure{Cel|Cela|Ceci} met en évidence que $H^u$ est un tore ``compact'', \deleted{\unsure{donc il est}} (de nature d'ailleurs assez particulière !), et aussi $U^0$ comme sous-tore de ce dernier\ldots.

NB. Il convient de déterminer $\varepsilon:G\longrightarrow\mathrm{Gm}$ en termes \inserted{de} \deleted{des} \unsure{leurs|les|tores} normes sur les $L_i$, et $j:\mathrm{Gm}\longrightarrow G$ \margin{$=\ldots$}
% archive sheet number: 85

